\section{Broader Impacts}

\paragraph{Improving Software System Reliability:} A key element of
this proposal's approach is to focus on realistically deployable
techniques.  All methods developed will aim for seamless integration
with the AFL++ fuzzing system, and thus immediate practical
applicability to real-world problems.  Differential testing is widely
used in various aspects of critical systems infrastructure, and
improvements in the ability to detect bugs using differential methods
are likely to contribute to software reliability.  Additionally, as
noted in the problem statement, it seems likely that differential
approaches are both made more practical by, and essential for the
reliability of, LLM and agentic coding efforts.  It is likely that
even future workflows with few humans in the loop will make use of
differential techniques, and given that fuzzing driven by random
mutations is likely an orthogonal source of bugs to LLM-based static
analysis of code, such methods will be more important for finding bugs
made by LLMs that they cannot reliably detect by only their own
``native'' methods.

\paragraph{Education and Outreach:}
The proposed research yields several opportunities for enhancing CS
education, recruiting new CS majors, and retaining CS students.  PI Groce will work with the NAU Student ACM Chapter to
present a series of ``excursions in testing'' that introduce automated
testing to students, using differential approaches to find bugs in real world code,
including code from media player libraries.  The
work of Guzdial \cite{Guzdial} has shown that media computation is an
effective way to both recruit and retain students in computer science. Groce also teaches
classes (undergraduate and graduate, including a required course in
NAU's Software Engineering major) on automated testing that can
present project results.  Groce's prior
engagement in teaching and mentoring high school students (including
as summer students who have published their work in major venues) and working with
local technically-oriented 4H chapters, including in rural communities
(such as Williams, AZ) surrounding Flagstaff will also allow
preparation and presentation of versions of these excursions tailored
to high school students' level of knowledge.